\documentclass[11pt,a4paper]{article}

\usepackage[utf8]{inputenc}
\usepackage[T1]{fontenc}
\usepackage[french]{babel}
\usepackage[a4paper,margin=2.1cm]{geometry}
\usepackage{amsmath,amssymb,amsthm}
\usepackage{enumitem}
\usepackage{tabularx}
\usepackage{array}
\usepackage{xcolor}
\usepackage{lmodern}
\usepackage{fancyhdr}
\usepackage{lastpage}
\usepackage{hyperref}
\usepackage{multicol}
\usepackage{titlesec}
\usepackage{mathrsfs}
\usepackage{tcolorbox}

\hypersetup{
    colorlinks=true,
    linkcolor=black,
    urlcolor=blue
}

\pagestyle{fancy}
\fancyhf{}
\lhead{Terminale spécialité -- Mathématiques}
\rhead{Corrigé détaillé du DM -- Analyse de fonctions}
\cfoot{\thepage/\pageref{LastPage}}

\titleformat{\section}{\large\bfseries}{\thesection.}{0.5em}{}
\titleformat{\subsection}{\normalsize\bfseries}{\thesubsection.}{0.5em}{}

\setlist[itemize]{leftmargin=1.8em}
\setlist[enumerate]{leftmargin=2em}

\newcolumntype{Y}{>{\raggedright\arraybackslash}X}

\begin{document}

\begin{center}
    {\LARGE \textbf{Corrigé détaillé du Devoir Maison -- Mathématiques}}\\[0.3em]
    {\Large \textbf{Terminale spécialité}}\\[0.5em]
    {\large \textbf{Thème : Analyse de fonctions}}\\[0.8em]
\end{center}



\section*{Corrigé du Problème 1 -- Étude complète d'une fonction logarithmique \hfill /25}

On considère la fonction \(f\) définie sur \(]0,+\infty[\) par
\[
f(x)=x-1-\ln(x).
\]

\subsection*{1. Premières observations}

\begin{enumerate}
    \item La fonction logarithme népérien est définie uniquement pour \(x>0\). Donc
    \[
    D_f=]0,+\infty[.
    \]

    \item On calcule :
    \[
    f(1)=1-1-\ln(1)=0-0=0.
    \]

    \item Lorsque \(x\to 0^+\), on sait que
    \[
    \ln(x)\to -\infty.
    \]
    Ainsi
    \[
    -\ln(x)\to +\infty.
    \]
    Comme \(x-1\to -1\), on obtient
    \[
    f(x)=x-1-\ln(x)\to +\infty.
    \]

    \item Lorsque \(x\to +\infty\), on sait que \(x\) domine très largement \(\ln(x)\), donc
    \[
    x-1-\ln(x)\to +\infty.
    \]
    Ainsi
    \[
    \lim_{x\to+\infty}f(x)=+\infty.
    \]
\end{enumerate}

\subsection*{2. Variations de la fonction}

\begin{enumerate}
    \item Pour \(x>0\),
    \[
    f'(x)=1-\frac{1}{x}=\frac{x-1}{x}.
    \]

    \item Sur \(]0,+\infty[\), on a toujours \(x>0\), donc le signe de \(f'(x)\) est celui de \(x-1\).

    Ainsi :
    \[
    f'(x)<0 \quad \text{si } 0<x<1,
    \]
    \[
    f'(1)=0,
    \]
    \[
    f'(x)>0 \quad \text{si } x>1.
    \]

    \item On en déduit que \(f\) est décroissante sur \(]0,1]\), puis croissante sur \([1,+\infty[\).

    Comme \(f(1)=0\), on obtient le tableau de variations :

    \[
    \begin{array}{c|ccc}
    x & 0^+ & 1 & +\infty \\
    \hline
    f'(x) & - & 0 & + \\
    \hline
    f(x) & +\infty & \searrow & \\
         &         & 0 & \nearrow +\infty
    \end{array}
    \]

    \item Le minimum de \(f\) sur \(]0,+\infty[\) est atteint en \(x=1\) et vaut \(0\). Donc, pour tout \(x>0\),
    \[
    f(x)\geq 0.
    \]
\end{enumerate}

\subsection*{3. Étude de la convexité}

\begin{enumerate}
    \item Pour \(x>0\),
    \[
    f''(x)=\frac{1}{x^2}.
    \]

    \item Comme \(x^2>0\) pour tout \(x>0\), on a
    \[
    f''(x)>0 \quad \text{pour tout } x>0.
    \]

    \item La fonction \(f\) est donc convexe sur \(]0,+\infty[\).

    \item Graphiquement, cela signifie que la courbe de \(f\) est tournée vers le haut et qu'elle est située au-dessus de chacune de ses tangentes.
\end{enumerate}

\subsection*{4. Tangente et comparaison locale}

\begin{enumerate}
    \item Une équation de la tangente en \(x=1\) est
    \[
    y=f'(1)(x-1)+f(1).
    \]
    Or
    \[
    f'(1)=0 \quad \text{et} \quad f(1)=0.
    \]
    Donc l'équation de la tangente est
    \[
    y=0.
    \]

    \item Étudier la position relative de la courbe et de cette tangente revient à étudier le signe de
    \[
    f(x)-0=f(x).
    \]
    Or on a montré que, pour tout \(x>0\),
    \[
    f(x)\geq 0.
    \]
    La courbe est donc au-dessus de sa tangente en \(x=1\), avec contact au point \((1,0)\).

    \item Comme
    \[
    f(x)=x-1-\ln(x)\geq 0,
    \]
    on obtient
    \[
    \ln(x)\leq x-1 \quad \text{pour tout } x>0.
    \]
\end{enumerate}

\subsection*{5. Regard critique}

\begin{enumerate}
    \item L'affirmation \og comme \(f(1)=0\), la fonction change forcément de signe en \(1\) \fg\ est fausse.

    En effet, le fait qu'une fonction s'annule en un point ne suffit pas à garantir un changement de signe. Ici, au contraire, on a montré que
    \[
    f(x)\geq 0 \quad \text{pour tout } x>0,
    \]
    donc la fonction ne change pas de signe en \(1\). Elle y atteint simplement un minimum nul.

    \item Ce problème montre géométriquement que la courbe du logarithme est en dessous de la droite d'équation
    \[
    y=x-1.
    \]
    Comme cette droite intervient via
    \[
    \ln(x)\leq x-1,
    \]
    on comprend que certaines tangentes ou droites d'appui permettent d'encadrer la fonction logarithme.
\end{enumerate}

\newpage

\section*{Corrigé du Problème 2 -- Une fonction rationnelle-exponentielle pour modéliser un rendement \hfill /29}

On considère la fonction \(g\) définie sur \([0,+\infty[\) par
\[
g(x)=\frac{x}{x+1}e^{-x}.
\]

\subsection*{1. Étude préliminaire}

\begin{enumerate}
    \item Pour \(x\geq 0\), on a \(x+1\neq 0\). Donc la fonction est bien définie sur
    \[
    [0,+\infty[.
    \]

    \item En \(x=0\),
    \[
    g(0)=\frac{0}{0+1}e^0=0.
    \]

    \item Lorsque \(x\to +\infty\), on a
    \[
    \frac{x}{x+1}\to 1
    \quad \text{et} \quad
    e^{-x}\to 0.
    \]
    Donc
    \[
    g(x)\to 0.
    \]

    \item Graphiquement, la courbe part de l'origine, monte d'abord, puis finit par redescendre vers \(0\). L'axe des abscisses est asymptote horizontale à droite.
\end{enumerate}

\subsection*{2. Dérivation et variations}

\begin{enumerate}
    \item On dérive par produit :
    \[
    g(x)=\left(\frac{x}{x+1}\right)e^{-x}.
    \]
    Or
    \[
    \left(\frac{x}{x+1}\right)'=\frac{(x+1)-x}{(x+1)^2}=\frac{1}{(x+1)^2}.
    \]
    Et
    \[
    (e^{-x})'=-e^{-x}.
    \]
    Donc
    \[
    g'(x)=\frac{1}{(x+1)^2}e^{-x}+\frac{x}{x+1}(-e^{-x}).
    \]
    On factorise par \(e^{-x}\) :
    \[
    g'(x)=e^{-x}\left(\frac{1}{(x+1)^2}-\frac{x}{x+1}\right).
    \]
    On met au même dénominateur :
    \[
    \frac{1}{(x+1)^2}-\frac{x(x+1)}{(x+1)^2}
    =\frac{1-x(x+1)}{(x+1)^2}.
    \]
    Donc
    \[
    g'(x)=\frac{e^{-x}}{(x+1)^2}\bigl(1-x-x^2\bigr).
    \]

    \item Pour \(x\geq 0\), on a
    \[
    e^{-x}>0 \quad \text{et} \quad (x+1)^2>0.
    \]
    Le signe de \(g'(x)\) est donc celui de
    \[
    1-x-x^2.
    \]
    Résolvons
    \[
    1-x-x^2=0 \iff x^2+x-1=0.
    \]
    Les racines sont
    \[
    x=\frac{-1\pm \sqrt{5}}{2}.
    \]
    La seule racine positive est
    \[
    \alpha=\frac{-1+\sqrt{5}}{2}.
    \]
    Comme le trinôme \(1-x-x^2\) est positif entre ses racines et négatif à l'extérieur, sur \([0,+\infty[\) on a :
    \[
    g'(x)>0 \quad \text{sur } [0,\alpha[,
    \]
    \[
    g'(\alpha)=0,
    \]
    \[
    g'(x)<0 \quad \text{sur } ]\alpha,+\infty[.
    \]

    \item Donc \(g\) est croissante sur \([0,\alpha]\), puis décroissante sur \([\alpha,+\infty[\).

    \item L'équation
    \[
    1-x-x^2=0
    \]
    admet bien une unique solution positive :
    \[
    \alpha=\frac{-1+\sqrt{5}}{2}.
    \]

    \item On en déduit que \(g\) admet un maximum sur \([0,+\infty[\), atteint en \(x=\alpha\).
\end{enumerate}

\subsection*{3. Exploitation du maximum}

\begin{enumerate}
    \item La valeur exacte est
    \[
    \alpha=\frac{-1+\sqrt{5}}{2}.
    \]

    \item Une valeur approchée au centième est
    \[
    \alpha\approx \frac{-1+2{,}236}{2}\approx 0{,}62.
    \]

    \item Calculons une valeur approchée de \(g(\alpha)\).

    Avec \(\alpha\approx 0{,}62\),
    \[
    \frac{\alpha}{\alpha+1}\approx \frac{0{,}62}{1{,}62}\approx 0{,}383,
    \]
    et
    \[
    e^{-0{,}62}\approx 0{,}538.
    \]
    Donc
    \[
    g(\alpha)\approx 0{,}383\times 0{,}538\approx 0{,}206.
    \]
    Ainsi
    \[
    g(\alpha)\approx 0{,}21.
    \]

    \item Dans le modèle, \(\alpha\) représente la valeur du paramètre \(x\) pour laquelle le rendement est maximal.
\end{enumerate}

\subsection*{4. Une équation associée}

On considère l'équation
\[
g(x)=0{,}1
\]
sur \([0,+\infty[\).

\begin{enumerate}
    \item Comme \(g\) est continue, croissante puis décroissante, et que son maximum est supérieur à \(0{,}1\), une droite horizontale peut couper la courbe en zéro, une ou deux fois selon sa hauteur.

    \item Ici :
    \[
    g(0)=0<0{,}1, \qquad g(\alpha)\approx 0{,}21>0{,}1, \qquad \lim_{x\to+\infty}g(x)=0<0{,}1.
    \]
    La courbe monte au-dessus de \(0{,}1\), puis redescend en dessous. Il y a donc deux solutions.

    \item Cherchons des encadrements simples.

    \textbf{Première solution :}

    \[
    g(0{,}1)=\frac{0{,}1}{1{,}1}e^{-0{,}1}\approx 0{,}091\times 0{,}905\approx 0{,}082,
    \]
    \[
    g(0{,}2)=\frac{0{,}2}{1{,}2}e^{-0{,}2}\approx 0{,}167\times 0{,}819\approx 0{,}137.
    \]
    Donc une première solution appartient à
    \[
    [0{,}1;0{,}2].
    \]

    \textbf{Deuxième solution :}

    \[
    g(2)=\frac{2}{3}e^{-2}\approx 0{,}667\times 0{,}135\approx 0{,}090,
    \]
    \[
    g(1{,}5)=\frac{1{,}5}{2{,}5}e^{-1{,}5}\approx 0{,}6\times 0{,}223\approx 0{,}134.
    \]
    Donc une seconde solution appartient à
    \[
    [1{,}5;2].
    \]
\end{enumerate}

\subsection*{5. Prise de recul}

\begin{enumerate}
    \item La fonction \(\dfrac{x}{x+1}\) seule tend vers \(1\), alors que le facteur \(e^{-x}\) tend vers \(0\) très rapidement. Ce facteur exponentiel force la fonction entière à redescendre vers \(0\), ce qui change complètement le comportement global.

    \item Si on remplaçait \(e^{-x}\) par \(e^{-2x}\), la décroissance serait plus rapide. Le maximum serait sans doute atteint plus tôt, et les valeurs de la fonction deviendraient plus petites pour les grandes valeurs de \(x\).
\end{enumerate}

\newpage

\section*{Corrigé du Problème 3 -- Une équation fonctionnelle autour des symétries d'une courbe \hfill /30}

On considère la fonction \(f\) définie sur \(\mathbb{R}\) par :
\[
f(x)=\frac{1}{1+e^{-x}}.
\]

\subsection*{1. Premières propriétés}

\begin{enumerate}
    \item Pour tout réel \(x\), on a \(e^{-x}>0\). Donc
    \[
    1+e^{-x}>0.
    \]
    Le dénominateur n'est donc jamais nul. La fonction \(f\) est définie sur \(\mathbb{R}\).

    \item En \(x=0\),
    \[
    f(0)=\frac{1}{1+e^0}=\frac{1}{2}.
    \]

    \item Lorsque \(x\to -\infty\), on a \(-x\to +\infty\), donc
    \[
    e^{-x}\to +\infty.
    \]
    Ainsi
    \[
    f(x)=\frac{1}{1+e^{-x}}\to 0.
    \]

    Lorsque \(x\to +\infty\), on a \(-x\to -\infty\), donc
    \[
    e^{-x}\to 0.
    \]
    Ainsi
    \[
    f(x)=\frac{1}{1+e^{-x}}\to 1.
    \]

    \item Graphiquement, la courbe admet deux asymptotes horizontales :
    \[
    y=0 \quad \text{à gauche}, \qquad y=1 \quad \text{à droite}.
    \]
\end{enumerate}

\subsection*{2. Variations de la fonction}

\begin{enumerate}
    \item On écrit
    \[
    f(x)=\bigl(1+e^{-x}\bigr)^{-1}.
    \]
    Donc
    \[
    f'(x)=-\bigl(1+e^{-x}\bigr)^{-2}\times(-e^{-x})
    =\frac{e^{-x}}{(1+e^{-x})^2}.
    \]

    \item Pour tout réel \(x\),
    \[
    e^{-x}>0 \quad \text{et} \quad (1+e^{-x})^2>0.
    \]
    Donc
    \[
    f'(x)>0 \quad \text{pour tout } x\in\mathbb{R}.
    \]

    \item La fonction \(f\) est donc strictement croissante sur \(\mathbb{R}\).

    Son tableau de variations est :
    \[
    \begin{array}{c|ccc}
    x & -\infty & & +\infty \\
    \hline
    f'(x) & & + & \\
    \hline
    f(x) & 0 & \nearrow & 1
    \end{array}
    \]
\end{enumerate}

\subsection*{3. Une équation fonctionnelle}

\begin{enumerate}
    \item Calculons
    \[
    f(-x)=\frac{1}{1+e^{-(-x)}}=\frac{1}{1+e^x}.
    \]

    \item Montrons que
    \[
    f(x)+f(-x)=1.
    \]
    On a
    \[
    f(x)=\frac{1}{1+e^{-x}}=\frac{e^x}{1+e^x}.
    \]
    Donc
    \[
    f(x)+f(-x)=\frac{e^x}{1+e^x}+\frac{1}{1+e^x}
    =\frac{e^x+1}{1+e^x}=1.
    \]

    \item C'est une équation fonctionnelle car elle relie les valeurs de la fonction en deux points liés, ici \(x\) et \(-x\).
\end{enumerate}

\subsection*{4. Interprétation graphique}

\begin{enumerate}
    \item Si \(M(x,f(x))\) appartient à la courbe, alors
    \[
    f(-x)=1-f(x).
    \]
    Donc le point
    \[
    M'(-x,1-f(x))=(-x,f(-x))
    \]
    appartient aussi à la courbe.

    \item Le milieu du segment \([MM']\) a pour coordonnées :
    \[
    \left(\frac{x+(-x)}{2},\frac{f(x)+1-f(x)}{2}\right)
    =
    \left(0,\frac12\right).
    \]
    Ce milieu est fixe. La courbe admet donc un centre de symétrie en
    \[
    \left(0,\frac12\right).
    \]

    \item Pour \(x=0\), on retrouve le point
    \[
    \left(0,\frac12\right),
    \]
    ce qui confirme bien le résultat.
\end{enumerate}

\subsection*{5. Une autre écriture utile}

\begin{enumerate}
    \item On calcule :
    \[
    1-f(x)=1-\frac{1}{1+e^{-x}}
    =\frac{1+e^{-x}-1}{1+e^{-x}}
    =\frac{e^{-x}}{1+e^{-x}}.
    \]

    \item Alors
    \[
    f(x)\bigl(1-f(x)\bigr)
    =\frac{1}{1+e^{-x}}\times \frac{e^{-x}}{1+e^{-x}}
    =\frac{e^{-x}}{(1+e^{-x})^2}.
    \]
    Or on a montré que
    \[
    f'(x)=\frac{e^{-x}}{(1+e^{-x})^2}.
    \]
    Donc
    \[
    f'(x)=f(x)\bigl(1-f(x)\bigr).
    \]

    \item Comme \(0<f(x)<1\), les deux facteurs \(f(x)\) et \(1-f(x)\) sont positifs, donc
    \[
    f'(x)>0.
    \]
    Cette écriture permet donc de comprendre simplement pourquoi la fonction est croissante.
\end{enumerate}

\subsection*{6. Question de recul}

\begin{enumerate}
    \item L'affirmation \og comme \(f(x)+f(-x)=1\), la fonction \(f\) est paire \fg\ est fausse.

    Une fonction paire vérifie
    \[
    f(-x)=f(x).
    \]
    Ici, on a au contraire
    \[
    f(-x)=1-f(x),
    \]
    ce qui est très différent.

    \item Il faut distinguer :
    \begin{itemize}
        \item la symétrie par rapport à l'axe des ordonnées, qui correspond à une fonction paire ;
        \item la symétrie centrale, qui signifie ici que la courbe est inchangée par une rotation de \(180^\circ\) autour du point
        \[
        \left(0,\frac12\right).
        \]
    \end{itemize}
\end{enumerate}

\subsection*{Conclusion}

La fonction
\[
f(x)=\frac{1}{1+e^{-x}}
\]
est définie et strictement croissante sur \(\mathbb{R}\), avec
\[
\lim_{x\to-\infty}f(x)=0
\qquad \text{et} \qquad
\lim_{x\to+\infty}f(x)=1.
\]
Elle vérifie l'équation fonctionnelle remarquable
\[
f(x)+f(-x)=1,
\]
ce qui montre que sa courbe admet un centre de symétrie en
\[
\left(0,\frac12\right).
\]
Enfin, l'identité
\[
f'(x)=f(x)\bigl(1-f(x)\bigr)
\]
fournit une autre lecture très intéressante de sa croissance.

\end{document}